\chapter*{Course Map}
\phantomsection%
\addcontentsline{toc}{chapter}{Course Map}

The course progresses from data organization to algorithm selection. Five stages build core skills; the final project compares algorithms with evidence.

\begin{mlfigure}
\begin{tikzpicture}[
  node distance=0.48cm and 0.72cm,
  box/.style={draw=ADSBlue,fill=ADSPaleBlue,rounded corners=2mm,text width=5.15cm,minimum height=0.92cm,align=left,inner xsep=4mm,font=\small\color{ADSNavy}},
  arr/.style={-{Stealth[length=2.2mm]},thick,draw=ADSTeal}
]
\node[box] (a) {\textbf{\color{ADSBlue}1. Foundations}\\[-0.15em]{\scriptsize Arrays and memory}};
\node[box,right=of a] (b) {\textbf{\color{ADSBlue}2. Linear structures}\\[-0.15em]{\scriptsize Lists, stacks, and queues}};
\node[box,below=of b] (c) {\textbf{\color{ADSBlue}3. Hierarchies}\\[-0.15em]{\scriptsize Trees, BSTs, and heaps}};
\node[box,left=of c] (d) {\textbf{\color{ADSBlue}4. Networks}\\[-0.15em]{\scriptsize Graph representation and traversal}};
\node[box,below=of d] (e) {\textbf{\color{ADSBlue}5. Core algorithms}\\[-0.15em]{\scriptsize Searching, sorting, and hashing}};
\node[box,right=of e] (f) {\textbf{\color{ADSBlue}6. Synthesis}\\[-0.15em]{\scriptsize Evidence-based comparison}};
\draw[arr] (a) -- (b);
\draw[arr] (b) -- (c);
\draw[arr] (c) -- (d);
\draw[arr] (d) -- (e);
\draw[arr] (e) -- (f);
\end{tikzpicture}
\end{mlfigure}

\begin{mltable}
\begin{tabularx}{\linewidth}{@{}>{\bfseries\leavevmode\color{ADSNavy}}p{0.08\linewidth}p{0.25\linewidth}X@{}}
\toprule
\rowcolor{ADSPaleBlue!92}
Stage & Focus & Key question \\
\midrule
1 & Arrays and memory & How is data stored and accessed? \\
2 & Lists, stacks, and queues & How do links and access rules organize data? \\
3 & Trees, BSTs, and heaps & How does hierarchy improve searching? \\
4 & Graphs, BFS, and DFS & How are relationships represented and traversed? \\
5 & Searching, sorting, and hashing & How does algorithm choice affect performance? \\
6 & Comparison project & How can algorithms be compared fairly? \\
\bottomrule
\end{tabularx}
\end{mltable}

\section*{Complexity vocabulary used across the course}
The table summarizes eight common Big-O classes, their representative growth functions, and typical algorithm examples.

\begin{mltable}
\begin{tabularx}{\linewidth}{@{}>{\bfseries\leavevmode\color{ADSNavy}}p{0.11\linewidth}p{0.20\linewidth}p{0.18\linewidth}X@{}}
\toprule
\rowcolor{ADSPaleBlue!92}
Class & Figure label & Representative function & Typical examples \\
\midrule
$O(1)$ & Constant & $f(n)=1$ & Array access; hash computation. \\
$O(\log n)$ & Logarithmic & $f(n)=\log_2 n$ & Binary search; balanced-BST lookup. \\
$O(n)$ & Linear & $f(n)=n$ & Linear search; linked-list traversal. \\
$O(n\log n)$ & Log-linear & $f(n)=n\log_2 n$ & Merge sort; heap sort. \\
$O(n^2)$ & Quadratic & $f(n)=n^2$ & Bubble sort; selection sort. \\
$O(n^3)$ & Cubic & $f(n)=n^3$ & Triple-nested loops; Floyd--Warshall. \\
$O(2^n)$ & Exponential & $f(n)=2^n$ & Subset enumeration; exhaustive include/exclude search. \\
$O(n!)$ & Factorial & $f(n)=n!$ & Permutation enumeration; brute-force traveling salesperson search. \\
\bottomrule
\end{tabularx}
\end{mltable}

\begin{mlfigure}
\centering
\resizebox{0.98\linewidth}{!}{%
\begin{tikzpicture}[
  line cap=round,
  line join=round,
  curve/.style={line width=1.65pt,-{Stealth[length=2.2mm,width=1.6mm]}},
  directlabel/.style={font=\footnotesize\bfseries,inner sep=0.8pt}
]
% Open axes keep the composition uncluttered and emphasize continued growth.
\draw[-{Stealth[length=2.9mm]},line width=1.05pt,draw=ADSMuted!88!black]
  (0,0) -- (12.1,0);
\draw[-{Stealth[length=2.9mm]},line width=1.05pt,draw=ADSMuted!88!black]
  (0,0) -- (0,8.1);
\node[font=\small\bfseries\color{ADSNavy},anchor=north] at (6.05,-0.43) {Input Size};
\node[font=\small\bfseries\color{ADSNavy},rotate=90,anchor=south] at (-0.64,4.05) {Computations};

% Every curve begins at the same conceptual lower-left point.
\coordinate (start) at (0.38,0.42);

% Constant and efficient growth.
\draw[curve,draw=ADSBlue] (start) -- (8.72,0.42);
\draw[curve,draw=green!48!black]
  (start) .. controls (1.02,1.25) and (3.22,1.48) .. (8.72,1.64);

% Linear, log-linear, and polynomial growth.
\draw[curve,draw=yellow!55!orange!78!black]
  (start) -- (8.72,2.86);
\draw[curve,draw=ADSOrange]
  (start) .. controls (3.40,1.08) and (6.65,2.82) .. (8.28,4.06);
\draw[curve,draw=ADSRed]
  (start) .. controls (3.95,0.53) and (7.02,2.35) .. (7.88,5.62);
% Cubic is one specific polynomial class; it grows faster than quadratic
% but remains asymptotically below exponential growth.
\draw[curve,draw=ADSRed!55!black]
  (start) .. controls (3.45,0.44) and (6.05,1.48) .. (7.18,7.20);

% Combinatorial growth.
\draw[curve,draw=magenta!82!red]
  (start) .. controls (2.22,0.41) and (4.47,0.89) .. (5.12,7.45);
\draw[curve,draw=purple!76!black]
  (start) .. controls (1.34,0.40) and (2.49,0.69) .. (2.72,7.47);

% Labels sit beyond curve endpoints so no text obscures a line.
\node[directlabel,text=purple!76!black,anchor=west] at (0.68,7.62)
  {Factorial -- $O(n!)$};
\node[directlabel,text=magenta!82!red,anchor=west] at (3.02,7.04)
  {Exponential -- $O(2^n)$};
\node[directlabel,text=ADSRed!55!black,anchor=west] at (7.34,7.20)
  {Cubic -- $O(n^3)$};
\node[directlabel,text=ADSRed,anchor=west] at (8.10,5.62)
  {Quadratic -- $O(n^2)$};
\node[directlabel,text=ADSOrange,anchor=west] at (8.50,4.06)
  {Log-linear -- $O(n\log n)$};
\node[directlabel,text=yellow!55!orange!78!black,anchor=west] at (8.94,2.86)
  {Linear -- $O(n)$};
\node[directlabel,text=green!48!black,anchor=west] at (8.94,1.64)
  {Logarithmic -- $O(\log n)$};
\node[directlabel,text=ADSBlue,anchor=west] at (8.94,0.42)
  {Constant -- $O(1)$};
\end{tikzpicture}%
}
\end{mlfigure}
